% F081: precise mathematical teaching diagram.
\begin{tikzpicture}[x=1cm,y=1cm,>=stealth,font=\small]
\draw[rounded corners=3pt,gray!40] (0,0) rectangle (14.2,8.5);
\begin{scope}[yshift=1cm]
\node at (7.1,6.9) {易错五：两个零点之间缺了一点};
\begin{scope}[shift={(7.1,3.9)},x=2.2cm,y=1.1cm]
\draw[->,gray!65] (-2.45,0)--(2.45,0) node[right] {$x$};
\draw[->,gray!65] (0,-2)--(0,2) node[above] {$y$};
\draw[blue!75!black,line width=1.15pt,,domain=-2:-0.002,samples=160] plot (\x,{\x+1});
\draw[blue!75!black,line width=1.15pt,,domain=0.002:2,samples=160] plot (\x,{\x-1});
\draw[blue!75!black,fill=white,line width=1pt] (0,1) circle (2.3pt) (0,-1) circle (2.3pt);
\fill[blue!75!black] (-1,0) circle (1.8pt);\node[below] at (-1,0) {$-1$};
\fill[blue!75!black] (1,0) circle (1.8pt);\node[above] at (1,0) {$1$};
\end{scope}
\node[] at (7.1,1.3) {$x=0$ 未定义，不能把两段拼成一个适用区间。};
\end{scope}
\node[text width=13cm,align=center] at (7.1,.9) {两段都满足 $f^{\prime}=1$；连完整区间都没有，不能跨缺口使用罗尔定理。};
\end{tikzpicture}
